\documentclass[11pt]{article}

\usepackage{amsmath, amssymb, amsthm}
\usepackage{mathtools}
\usepackage{enumitem}
\usepackage{hyperref}
\usepackage[margin=1.2in]{geometry}

% Fonts (uncomment if installed):
% \usepackage{fouriernc}
% \usepackage[T1]{fontenc}

% Theorem-like environments
\theoremstyle{definition}
\newtheorem*{solution}{Solution}
\newtheorem*{claim}{Claim}

% Common abbreviations
\newcommand{\N}{\mathbb{N}}
\renewcommand{\P}{\mathbb{P}}
\newcommand{\Z}{\mathbb{Z}}
\newcommand{\Q}{\mathbb{Q}}
\newcommand{\R}{\mathbb{R}}
\newcommand{\C}{\mathbb{C}}
\newcommand{\F}{\mathbb{F}}

% Formal power series rings
\newcommand{\Cfps}{\mathbb{C}[[x]]}
\newcommand{\Rfps}{R[[x]]}

% Generating function coefficient extraction: \coeff{x^n}{F(x)}
\newcommand{\coeff}[2]{[{#1}]\,{#2}}

% Combinatorics
\newcommand{\falling}[2]{{#1}^{\underline{#2}}}
\newcommand{\rising}[2]{{#1}^{\overline{#2}}}
\newcommand{\stirlingI}[2]{\genfrac{[}{]}{0pt}{}{#1}{#2}}
\newcommand{\stirlingII}[2]{\genfrac{\{}{\}}{0pt}{}{#1}{#2}}

% Multiset coefficient ((n choose m))
\newcommand{\multiset}[2]{\left(\!\binom{#1}{#2}\!\right)}

% Convenience
\DeclareMathOperator{\ord}{ord}
\newcommand{\eps}{\varepsilon}
\newcommand{\ds}{\displaystyle}

% Problem wrapper: \problem{N} produces "Problem 1.N."
\newcounter{probnum}
\newcommand{\problem}[1]{%
  \setcounter{probnum}{#1}%
  \medskip
  \noindent\textbf{Problem 1.#1.}\quad
}

\title{Stanley --- Enumerative Combinatorics, Vol.\ 1\\
       \large Chapter 1 Solutions}
\author{}
\date{}

\begin{document}
\maketitle

\problem{1} \textit{[1--]}
Let $S$ and $T$ be disjoint one-element sets.
Find the number of elements of their union $S \cup T$.

\begin{solution}

This is clear.  2.

\end{solution}

\problem{2} \textit{[1+]}
Find as simple a solution as possible.

\begin{enumerate}[label=(\alph*)]
  \item How many subsets of the set $[10] = \{1,2,\ldots,10\}$ contain at least one odd integer?
    \begin{solution}
    This is total sets minus total with only even $= 2^{10} - 2^5 = 1024 - 32 = 992.$
    \end{solution}

  \item In how many ways can seven people be seated in a circle if two arrangements are
    considered the same whenever each person has the same neighbors (not necessarily
    on the same side)?
    \begin{solution}
    Assuming there are named seats, this means a chain of seven in either one order or the other (if A is to the left of B to the left of C to the left of D, then reversing so C is to the left of B is to the left of A means that D is must be to the left of C).  One person is in seat 1, so this is $7*2 = 14$.  If the seats are unnamed, this 2.
    \end{solution}

  \item How many permutations $w:[6]\to[6]$ satisfy $w(1)\neq 2$?
    \begin{solution}
    Permutations with $w(1) = 2$ are $5!$.  Total perms is $6!$.  Difference is $6! - 5! = 600$.
    \end{solution}

  \item How many permutations of $[6]$ have exactly two cycles (i.e., find $c(6,2)$)?
    \begin{solution}
    $c(n,1) = (n-1)!$ in standard form, since it starts with the highest element.
    Using the recurrence $c(n,k) = (n-1)c(n-1,k) + c(n-1,k-1)$:
    \begin{align*}
      c(2,2) &= 1 \\
      c(3,2) &= 2c(2,2) + c(2,1) = 2 + 1 = 3 \\
      c(4,2) &= 3c(3,2) + c(3,1) = 9 + 2 = 11 \\
      c(5,2) &= 4c(4,2) + c(4,1) = 44 + 6 = 50 \\
      c(6,2) &= 5c(5,2) + c(5,1) = 250 + 24 = 274
    \end{align*}
    Confirmed by finding the coefficient of $t^2$ in $t(t+1)\cdots(t+5)$.
    \end{solution}

  \item How many partitions of $[6]$ have exactly three blocks (i.e., find $S(6,3)$)?
    \begin{solution}
    Using the recurrence $S(n,k) = kS(n-1,k) + S(n-1,k-1)$ (Prop.\ 1.9.3) or the
    inclusion-exclusion formula (1.94a):
    \begin{align*}
      S(6,3) &= \frac{1}{3!}\sum_{j=0}^{3}(-1)^{3-j}\binom{3}{j}j^6 \\
             &= \frac{1}{6}\left[-\binom{3}{0}0^6 + \binom{3}{1}1^6
                - \binom{3}{2}2^6 + \binom{3}{3}3^6\right] \\
             &= \frac{0 + 3 - 192 + 729}{6} = \frac{540}{6} = 90.
    \end{align*}
    Cross-check via the EGF (1.94b): $[x^6/6!](e^x-1)^3/3! = 6/48 \cdot 1/720$,
    giving $S(6,3) = 6!/48 \cdot (1/6) \cdot 6 = 90$. \checkmark
    \end{solution}


  \item There are four men and six women. Each man marries one of the women. In how
    many ways can this be done?
    \begin{solution}
    $6*5*4*3 =360$
    \end{solution}

  \item Ten people split up into five groups of two each. In how many ways can this be done?
    \begin{solution}
    Assuming groups are indistinct this is $\frac{10!}{ (2!)^5} =113,400 $
    \end{solution}

  \item How many compositions of 19 use only the parts 2 and 3?
    \begin{solution}
    % TODO: is there a cleaner approach?
    If we use $a$ threes and $b$ twos, then $3a+2b=19$, so $a$ must be odd.
    The number of compositions with $a$ threes and $b$ twos is $\binom{a+b}{a}$.
    \begin{center}
    \begin{tabular}{ccc}
      $a$ & $b$ & compositions \\ \hline
      1 & 8 & $\binom{9}{1} = 9$ \\
      3 & 5 & $\binom{8}{3} = 56$ \\
      5 & 2 & $\binom{7}{5} = 21$
    \end{tabular}
    \end{center}
    Total: $9 + 56 + 21 = 86$.
    \end{solution}

  \item In how many different ways can the letters of the word MISSISSIPPI be arranged
    if the four S's cannot appear consecutively?
    \begin{solution}
    MISSISSIPPI has 4 S's, 4 I's, 2 P's, 1 M\@. Strategy: subtract arrangements
    where all four S's appear consecutively.
    \begin{align*}
      \text{Total unrestricted:} &\quad \frac{11!}{4!\,4!\,2!} = 34650 \\
      \text{Violations (treat SSSS as one block):} &\quad \frac{8!}{4!\,2!} = 840 \\
      \text{Valid arrangements:} &\quad 34650 - 840 = 33810
    \end{align*}
    \end{solution}

  \item How many sequences $(a_1, a_2, \ldots, a_{12})$ are there consisting of four 0's and
    eight 1's, if no two consecutive terms are both 0's?
    \begin{solution}
    Place the eight 1's first; they create nine gaps (including the ends).
    Choose four of those gaps for the 0's:
    \[ \binom{9}{4} = 126. \]
    Equivalently, this is $[x^4](1+x)^9 = \binom{9}{4} = 126$.
    \end{solution}

  \item A box is filled with three blue socks, three red socks, and four chartreuse socks.
    Eight socks are pulled out, one at a time. In how many ways can this be done?
    (Socks of the same color are indistinguishable.)
    \begin{solution}
    Label the ten socks $1,\ldots,10$. The number of distinct ordered draws of all
    ten, treating same-color socks as identical, is
    \[ \frac{10!}{3!\,3!\,4!} = 4200. \]
    \end{solution}

  \item How many functions $f:[5]\to[5]$ are at most two-to-one, i.e., $\#f^{-1}(n)\leq 2$
    for all $n\in[5]$?
    \begin{solution}
    Model as labeled balls (domain) into labeled boxes (codomain), each box holding
    at most 2. The EGF for each box is $1 + x + \tfrac{x^2}{2!}$, so the count is
    \[ 5!\,[x^5]\!\left(1 + x + \tfrac{x^2}{2!}\right)^{\!5} = 2220. \]
    \end{solution}
\end{enumerate}

\problem{3}
Give combinatorial proofs of the following identities, where $x,y,n,a,b$ are
nonnegative integers.

\begin{enumerate}[label=(\alph*)]
  \item \textit{[2--]} $\displaystyle\sum_{k=0}^{n}\binom{x+k}{k} = \binom{x+n+1}{n}$
    \begin{solution}
    The right hand side indicates all paths of 0 (left) or 1 (right) of $x+n+1$
    with $n$ right choices.

    The left hand side, at index $k$, counts $n-k$ rights (1) before the first
    left, then $k$ rights after that.
    \end{solution}

  \item \textit{[1+]} $\displaystyle\sum_{k=0}^{n}k\binom{n}{k} = n\cdot 2^{n-1}$
    \begin{solution}
    Rewrite as $\sum_{k=0}^{n} \frac{k}{n}\binom{n}{k} = 2^{n-1}$. Further rewrite
    as $\sum_{k=0}^{n} \frac{k}{n}\binom{n}{k} = 2^{n-1}$ (note: $k=-1$ but that's
    0, so reindex).

    The left can be interpreted as the number of 0/1 combos of $n$ without selecting
    the first element; of all $\binom{n}{k}$ choices, $k/n$ of them will select the
    first element. With those eliminated, the total number of selections of $n$ 0/1s
    is $2^{n-1}$.
    \end{solution}

  \item \textit{[3]} $\displaystyle\sum_{k=0}^{n}\binom{2k}{k}\binom{2(n-k)}{n-k} = 4^n$
    \begin{solution}
    % TODO: apparently very hard; can be approached as NE/NW/SE/SW walks on a grid.
    TODO\@.
    \end{solution}

  \item \textit{[3--]}
    $\displaystyle\sum_{k=0}^{m}\binom{x}{b-k}\binom{y}{a-k}\binom{x+y+k}{k}
     = \binom{x+a}{a}\binom{y+b}{b}$,
    where $m = \min(a,b)$
    \begin{solution}
    TODO\@.
    \end{solution}

  \item \textit{[1]} $\displaystyle 2\binom{2n-1}{n} = \binom{2n}{n}$
    \begin{solution}
    Rewrite as $\binom{2n-1}{n} + \binom{2n-1}{n-1}$. Then $\binom{2n}{n}$ is the
    count of the $2n$-th element not in the choice (first term) plus in the choice
    (second term).
    \end{solution}

  \item \textit{[2--]} $\displaystyle\sum_{k=0}^{n}(-1)^{k}\binom{n}{k} = 0$, \quad $n\geq 1$
    \begin{solution}
    Sending the negative terms to the RHS, this is ``$n$ choose an even number''
    $=$ ``$n$ choose an odd number.'' This is clear if $n$ is odd (invert the
    choices, a bijection). If $n$ is even, you can assume by induction this is true
    at e.g.\ $n=2$, and then expanding $n$ by 2, you have the same number if you
    select 00 or 11 as selecting 01 or 10 with the two new elements.
    \end{solution}

  \item \textit{[2+]}
    $\displaystyle\sum_{k=0}^{n}x^k\binom{n}{k}
     = \sum_{j=0}^{n}\binom{n}{j}\binom{2n-j}{n}(x-1)^j$
    \begin{solution}
    TODO\@.

    Special case: $x=1$. Then this is choosing $n$ elements among $2n$:
    \begin{itemize}
      \item LHS: Split $2n$ into two, and choose $k$ from each half:
        $\binom{n}{k}\binom{n}{n-k} \cdot 1$.
      \item RHS: Only $j=0$ gets $(x-1)^j = 1$. So this is $\binom{n}{0}\binom{2n}{n}$.
    \end{itemize}
    These are equal. TODO: try with $x = 2^k$.
    \end{solution}

  \item \textit{[3--]}
    $\displaystyle\sum_{\substack{i,j,k\,\geq\,0\\i+j+k=n}}
       \binom{i+j}{i}\binom{j+k}{j}\binom{k+i}{k}
     = \sum_{r=0}^{n}\binom{2r}{r}$,
    \quad $i,j,k\in\N$
    \begin{solution}
    TODO\@.
    \end{solution}
\end{enumerate}

\problem{5} \textit{[2]*}
Show that
\[
  \sum_{n_1,\ldots,n_k\,\geq\,0} \min(n_1,\ldots,n_k)\, x_1^{n_1}\cdots x_k^{n_k}
  = \frac{x_1\cdots x_k}{(1-x_1)\cdots(1-x_k)(1-x_1 x_2\cdots x_k)}.
\]

\begin{solution}
Expand the right-hand side as
\[
  \frac{x_1 x_2 \cdots x_k}{(1-x_1 x_2 \cdots x_k)(1-x_1)\cdots(1-x_k)}
  = \bigl(x_1 x_2 \cdots x_k + (x_1 x_2 \cdots x_k)^2 + \cdots\bigr)
    (1+x_1+x_1^2+\cdots)\cdots(1+x_k+x_k^2+\cdots).
\]

First, we can consider only $n_1,\ldots,n_k > 0$, since setting any $n_i$ to zero
renders the left-hand term zero, and no term containing the factor $x_i^0$ exists on
the right side by virtue of the first factor including each $x_i$.

Both sides consist of exactly one term for each $(n_1,\ldots,n_k) \in \mathbb{N}^k$
with all entries positive: the left by naming each one in the sum, and the right by
collecting all terms into $x_1^{n_1} x_2^{n_2} \cdots x_k^{n_k}$. The only question
to resolve is the coefficient on the left-hand side.

For every $(n_1,\ldots,n_k)$, in the expanded form of the right-hand side (where all
coefficients are 1), the contribution from
$(x_1 x_2 \cdots x_k + (x_1 x_2 \cdots x_k)^2 + \cdots)$ will be from the power 1
to $\min(n_1,\ldots,n_k)$; the difference will be selected from the remaining factors.
There cannot be another selection from this first factor.

Therefore the coefficient on the left-hand side is $\min(n_1,\ldots,n_k)$, and the
equivalence is proven.
\end{solution}

\problem{7} \textit{[2]*}
Let
\[
  e^{x + x^2/2} = \sum_{n\geq 0} f(n)\,\frac{x^n}{n!}.
\]
Find a simple expression for $\displaystyle\sum_{i=0}^{n}(-1)^{n-i}\binom{n}{i}f(i)$.
(See equation (1.13).)

\begin{solution}
We have $\sum_{n=0}^{\infty} \frac{f(n)}{n!}x^n = e^{x+x^2/2}$.

What we need to do eventually is set $x$ to 1 and get each $x^n$ term of a generating
function as its own sum over $i$, with $f(i)$ getting an additional factor
$\binom{n}{i}(-1)^{n-i}$. We can't do this by changing $x$, and we can't square the
whole thing (that gives $f(i)f(j)$ terms), so we have to multiply it by an outside
function: $e^{-x}$.

This introduces exactly these terms, since combining two generating functions with
coefficients $a(i)$ and $b(n-i)$ for term $x^n$ yields $\binom{n}{i}$:
\[
  \sum_{n=0}^{\infty} \frac{x^n}{n!} \sum_{i=0}^{n} \binom{n}{i}f(i)(-1)^{n-i}
  = e^{x+x^2/2} \cdot e^{-x} = e^{x^2/2}.
\]
This has only even-degree terms, so the sum is zero for any odd $n$.

For even $n$, each term is $\dfrac{1}{2^n n!}x^{2n}$. Converting to EGF form,
$\dfrac{(2n)!}{2^n n!} \cdot \dfrac{x^{2n}}{(2n)!}$.

We can see that the coefficient $\dfrac{(2n)!}{2^n n!}$ is $\prod_{j=1}^{n}(2j-1)$.
\end{solution}

\problem{8} \hfill

\begin{enumerate}[label=(\alph*)]
  \item \textit{[2--]} Show that
    $\displaystyle\sum_{n\geq 0}\binom{2n}{n}x^n = \frac{1}{\sqrt{1-4x}}$.
    \begin{solution}
    By taking the Taylor series of the LHS, it's clear that $f^{(n)}(x)$ at $x=1$
    is $2^n \cdot (1 \cdot 3 \cdots (2n-1))$.

    There's an inductive proof that this product is $(n+1)(n+2)\cdots(2n)$: base
    case at $n=1$ is 2 on both sides. For the inductive case, divide by $2n$ to get
    the inductive case for the previous number.

    \textit{Lemma.}
    $\displaystyle 2^k \cdot \prod_{i=1}^{k-1}(2i+1) = \prod_{i=k+1}^{2k} i$.

    \textit{Proof of lemma.}
    \begin{itemize}
      \item Base case $k=1$: $2 = 2$. Done.
      \item We need to show $2^{k} \cdot \prod_{i=1}^{k}(2i+1) = \prod_{i=k+2}^{2k+2} i$
        using the inductive fact.
      \item This is the same as showing
        $2(2k+1)2^{k} \cdot \prod_{i=1}^{k-1}(2i+1) = (2k+2)(2k+1)\prod_{i=k+2}^{2k} i$.
      \item This is the same as showing
        $2(2k+1)2^{k} \cdot \prod_{i=1}^{k-1}(2i+1) = 2(2k+1)(2k+1)\prod_{i=k+1}^{2k} i$.
        Divide out $2(2k+1)$ to get to the inductive case.
    \end{itemize}

    Therefore $f^{(n)}x = 2^n\cdot \prod_{i=1}^{n-1}(2i+1)$, and
    \[
      \frac{f^{(n)}x}{n!}
        = 2^n\cdot \prod_{i=1}^{n-1}(2i+1) % subscript unclear in original
        = \frac{(n+1)(n+2)\cdots(2n)}{n!} = \binom{2n}{n}.
    \]
    Further, $\left((1-4x)^{-\frac{1}{2}}\right)^2 = \frac{1}{1-4x} = 1+4x+16x^2+\cdots$,
    so looking at each coefficient,
    $\sum_{k=0}^{n}\binom{2k}{k}\binom{2(n-k)}{n-k} = 4^n$.
    \end{solution}

  \item \textit{[2--]} Find
    $\displaystyle\sum_{n\geq 0}\binom{2n-1}{n}x^n$.
    \begin{solution}
    Since $\binom{2n-1}{n} = \frac{1}{2}\binom{2n}{n}$ for $n\geq 1$, subtract out
    the terms at $n=0$ to obtain the sum as
    $\dfrac{1}{2} + \dfrac{1}{2}(1-4x)^{-\frac{1}{2}}$.
    \end{solution}
\end{enumerate}

\problem{12} \textit{[2+]*}
Choose $n$ points on the circumference of a circle in ``general position.''
Draw all $\binom{n}{2}$ chords connecting two of the points.
(``General position'' means that no three of these chords intersect in a point.)
Into how many regions will the interior of the circle be divided?
Try to give an elegant proof avoiding induction, finite differences, generating functions,
summations, etc.

\begin{solution}
The solution is $\binom{n}{0} + \binom{n}{2} + \binom{n}{4}$.

Imagine drawing each of these chords and seeing where existing regions are split
(adding one more region).
\begin{itemize}
  \item $\binom{n}{0}$: There is one initial region inside the circle.
  \item $\binom{n}{2}$: Every time a chord touches the opposite vertex on the
    circle, a region is created. There are $\binom{n}{2}$ chords.
  \item $\binom{n}{4}$: Every time a chord crosses another chord the opposite
    vertex on the circle, a region is created. Every set of four vertices has
    one crossing point.
\end{itemize}
\end{solution}

\problem{13} \textit{[2]}
Let $p$ be prime and $a\in\P$. Show combinatorially that $a^p - a$ is divisible
by $p$. (A combinatorial proof would consist of exhibiting a set $S$ with $a^p - a$
elements and a partition of $S$ into pairwise disjoint subsets, each with $p$
elements.)

\begin{solution}
Imagine a loop of $p$ beads with $a$ different colors. You can specify it in $a^p$
ways starting from some index 1. However, there are only $\frac{a^p - a}{p}$ unique
loops irrespective of starting index: once you subtract out the uniform color bead
sequences ($a$), each loop, when rotated, generates a unique sequence; there can be
no matches on rotation because length $p$ is prime. This $a^p - a$ is divisible
by $p$.
\end{solution}

\problem{14}
Let $p$ be a prime, and let $n = \sum a_i p^i$ and $m = \sum b_i p^i$ be the
$p$-ary expansions of the positive integers $m$ and $n$.

\begin{enumerate}[label=(\alph*)]
  \item \textit{[2+]} Show that
    \[ \binom{n}{m} \equiv \binom{a_0}{b_0}\binom{a_1}{b_1}\cdots \pmod{p}. \]
    \begin{solution}
    \end{solution}

  \item \textit{[3--]} Use (a) to determine when $\binom{n}{m}$ is odd. For what $n$
    is $\binom{n}{m}$ odd for all $0\leq m\leq n$? In general, how many coefficients
    of the polynomial $(1+x)^n$ are odd?
    \begin{solution}
    \end{solution}

  \item \textit{[2+]} It follows from (a), and is easy to show directly, that
    $\binom{pa}{pb} \equiv \binom{a}{b} \pmod{p}$. Give a combinatorial proof that
    in fact $\binom{pa}{pb} \equiv \binom{a}{b} \pmod{p^2}$.
    \begin{solution}
    \end{solution}

  \item \textit{[3--]} If $p\geq 5$, show in fact
    $\binom{pa}{pb} \equiv \binom{a}{b} \pmod{p^3}$. Is there a combinatorial proof?
    \begin{solution}
    \end{solution}

  \item \textit{[3--]} Give a simple description of the largest power of $p$ dividing
    $\binom{n}{m}$.
    \begin{solution}
    \end{solution}
\end{enumerate}

\problem{15} \hfill

\begin{enumerate}[label=(\alph*)]
  \item \textit{[2]} How many coefficients of the polynomial $(1+x+x^2)^n$ are not
    divisible by~3?
    \begin{solution}
    \end{solution}

  \item \textit{[3--]} How many coefficients of the polynomial $(1+x+x^2)^n$ are odd?
    \begin{solution}
    \end{solution}

  \item \textit{[2+]} How many coefficients of the polynomial
    $\prod_{1\leq i<j\leq n}(x_i+x_j)$ are odd?
    \begin{solution}
    \end{solution}
\end{enumerate}

\problem{17}
Let $m, n \in \N$.

\begin{enumerate}[label=(\alph*)]
  \item \textit{[1+]*} Prove the identity
    \[ \multiset{n}{m} = \multiset{m+1}{n-1}. \]
    \begin{solution}
    \end{solution}

  \item \textit{[2--]} Give a combinatorial proof.
    \begin{solution}
    \end{solution}
\end{enumerate}

\problem{18} \textit{[2+]*}
Find a simple description of all $n\in\P$ with the following property: there exists
$k\in[n]$ such that $\binom{n}{k-1}$, $\binom{n}{k}$, $\binom{n}{k+1}$ are in
arithmetic progression.

\begin{solution}
$\binom{n}{k+1} - \binom{n}{k} = \binom{n}{k} - \binom{n}{k-1}$.
Refactoring: $(k+1)k + (n-k+1)(n-k) = 2(n-k+1)(k+1)$.
Refactoring: $(n-2k)^2 = n+2$.
So it happens whenever $n+2$ is a perfect square, i.e.\ $n = a^2 - 2$ for $a \geq 3$.
\end{solution}

\problem{20} \textit{[2]*}
How many $m\times n$ matrices of 0's and 1's are there such that every row and column
contains an even number of 1's? An odd number of 1's?

\begin{solution}
Proof idea: For an $n\times n$ matrix, choose a submatrix of $(n-1)\times(n-1)$
and set arbitrarily.

The subarray has an even or an odd total sum.

If odd, there are an odd number of odd total columns and an odd number of odd total
rows. Add one to the additional row or column wherever so. Add a one in the upper
right corner to make the whole matrix even.

If even, there are an even number of odd total columns and rows. Add one to the
additional column wherever so. Leave the corner 0.

In this way we produce an $n\times n$ matrix with even row and column totals.

If we want to produce odd totals, take the even matrix and invert the final row,
column, and corner.

Since the new row and column are determined by the submatrix, we can arbitrarily
choose the submatrix and no more, so the total number is at most $2^{(n-1)^2}$. But
since we can do this for any arbitrary $(n-1)\times(n-1)$, it's at least
$2^{(n-1)^2}$. So the answer is $2^{(n-1)^2}$.
\end{solution}

% Add further problems below as needed:
% \problem{N}
% \begin{solution}
% \end{solution}

\end{document}
